%% 1-17 表と ぼうグラフ（10ページ、p.172〜181）　第4部「表・グラフ・かたち」のはじめの章
%% 第1巻のきまり：分かち書き、漢字はすべてルビ、式は全角の＋－＝、数と単位のあいだは \,
%% 路線：⑤表。しらべる → かぞえる（正の字・／）→ 表 → ぼうグラフ。1目もりの大きさを決める。
%% はてな？の数：カレー 8、あげパン 7、ラーメン 5、やきそば 4（ぜんぶで 24まい）

\makeatletter
%% ---- 1-17 で使う部品（名前は \bkqp…） ----
% きゅう食の名前  \bkqpname{1〜4}   1 カレー  2 あげパン  3 ラーメン  4 やきそば
\newcommand\bkqpname[1]{\ifcase#1\or カレー\or あげパン\or ラーメン\or やきそば\fi}
\newcommand\bkqp@setcolor[1]{\ifcase#1\or\colorlet{bkqp@c}{bkOrange}\or\colorlet{bkqp@c}{bkStem}\or\colorlet{bkqp@c}{bkRed}\or\colorlet{bkqp@c}{bkGreen}\fi}
% きゅう食の小さな絵  \bkqpfood{(x,y)}{大きさ（1で はば約6.4mm）}{1〜4}
\newcommand\bkqpfood[3]{\begin{scope}[shift={#1}, scale=#2]
  \ifcase#3\or
    \filldraw[fill=white, draw=bkGray!70, line width=0.4pt] (0,0) ellipse (3.2 and 2.1);
    \filldraw[fill=white, draw=bkGray!45, line width=0.3pt] (-1.15,0.15) ellipse (1.55 and 1.2);
    \fill[bkOrange!70!bkStem] (1.15,-0.15) ellipse (1.6 and 1.2);
    \fill[bkOrange] (0.75,0.2) circle (0.33); \fill[bkYellow!55] (1.65,-0.55) circle (0.36);
  \or
    \filldraw[fill=bkYellow!60!bkStem, draw=bkStem, line width=0.4pt, rounded corners=1.3] (-3,-1.35) rectangle (3,1.35);
    \foreach \p in {(-2.1,0.5),(-1,-0.45),(0,0.55),(1.1,-0.35),(2.1,0.45),(-1.9,-0.6),(1.9,-0.7)} {\fill[white] \p circle (0.24);}
  \or
    \fill[bkYellow!45] (0,0.55) ellipse (3.1 and 0.95);
    \draw[bkYellow!75!black, line width=0.35pt] (-2.3,0.5) .. controls (-1.8,0.95) and (-1.3,0.1) .. (-0.8,0.5) .. controls (-0.3,0.95) and (0.2,0.1) .. (0.7,0.5);
    \filldraw[fill=white, draw=bkRed!70, line width=0.3pt] (1.75,0.65) circle (0.48);
    \fill[bkRed] (-3.1,0.55) arc (180:360:3.1 and 2.5) -- cycle;
    \draw[white, line width=0.35pt] (-2.2,-0.6) -- (2.2,-0.6);
  \or
    \filldraw[fill=white, draw=bkGray!70, line width=0.4pt] (0,0) ellipse (3.2 and 2.1);
    \foreach \yy in {-0.75,0,0.75} {\draw[bkStem!85, line width=0.55pt] (-2.1,\yy) .. controls (-1.4,\yy+0.45) and (-0.7,\yy-0.45) .. (0,\yy) .. controls (0.7,\yy+0.45) and (1.4,\yy-0.45) .. (2.1,\yy);}
    \fill[bkLeaf] (-0.9,0.35) circle (0.22) (0.6,-0.4) circle (0.22); \fill[bkRed!85] (1.3,0.75) circle (0.3);
  \fi
\end{scope}}
% こくばんに はる カード  \bkqpcard{(x,y)}{かたむき}{1〜4}
\newcommand\bkqpcard[3]{\begin{scope}[shift={#1}, rotate=#2]\bkqp@setcolor{#3}%
  \filldraw[fill=white, draw=bkqp@c, line width=1pt, rounded corners=0.8mm] (-9.8,-3.5) rectangle (9.8,3.5);
  \fill[bkqp@c] (0,3.5) circle (0.85);
  \bkqpfood{(-5.9,0)}{0.82}{#3}
  \node[anchor=west, font=\sffamily\bfseries\fontsize{8}{9}\selectfont, text=bkInk, inner sep=0pt] at (-2.3,0) {\bkqpname{#3}};
\end{scope}}
% 正の字  \bkqpsei[1字の大きさ mm]{数}   文の中や表の中で使う
\newcommand\bkqp@strokes[3]{\begin{scope}[shift={#1}]
  \draw[bkInk, line width=1pt, line cap=round] (0.1*#2,#2) -- (0.9*#2,#2);
  \ifnum#3>1 \draw[bkInk, line width=1pt, line cap=round] (0.5*#2,#2) -- (0.5*#2,0);\fi
  \ifnum#3>2 \draw[bkInk, line width=1pt, line cap=round] (0.5*#2,0.52*#2) -- (0.86*#2,0.52*#2);\fi
  \ifnum#3>3 \draw[bkInk, line width=1pt, line cap=round] (0.17*#2,0.6*#2) -- (0.17*#2,0);\fi
  \ifnum#3>4 \draw[bkInk, line width=1pt, line cap=round] (0,0) -- (#2,0);\fi
\end{scope}}
\newcommand\bkqpsei[2][3.8]{\tikz[x=1mm, y=1mm, baseline=0.2mm]{%
  \pgfmathtruncatemacro\bkqp@f{#2/5}\pgfmathtruncatemacro\bkqp@r{#2-5*\bkqp@f}%
  \path (0,-0.5) (0,#1+0.5);
  \ifnum\bkqp@f>0 \foreach \i in {1,...,\bkqp@f} {\bkqp@strokes{({(\i-1)*#1*1.45},0)}{#1}{5}}\fi
  \ifnum\bkqp@r>0 \bkqp@strokes{({\bkqp@f*#1*1.45},0)}{#1}{\bkqp@r}\fi}}
% グラフの目もり  \bkqpaxes{(x,y)}{はば}{目もりの数}{1目もりの高さ mm}{数字を書く間（目もり）}{1目もりの数（0なら0だけ書く）}{単位}
\newcommand\bkqpaxes[7]{\begin{scope}[shift={#1}]
  \foreach \k in {1,...,#3} {\pgfmathtruncatemacro\bkqp@m{mod(\k,#5)}
    \ifnum\bkqp@m=0 \draw[bkGray!60, line width=0.5pt] (0,\k*#4) -- (#2,\k*#4);
    \else\draw[bkLine, line width=0.45pt] (0,\k*#4) -- (#2,\k*#4);\fi}
  \ifnum#6>0
    \foreach \k in {0,#5,...,#3} {\pgfmathtruncatemacro\bkqp@v{\k*#6}
      \node[anchor=east, font=\sffamily\bfseries\fontsize{8.5}{9}\selectfont, text=bkInk, inner sep=0pt] at (-1.4,\k*#4) {\bkqp@v};}
  \else
    \node[anchor=east, font=\sffamily\bfseries\fontsize{8.5}{9}\selectfont, text=bkInk, inner sep=0pt] at (-1.4,0) {0};
  \fi
  \draw[bkInk, line width=0.9pt] (0,#3*#4+1.5) -- (0,0) -- (#2,0);
  \node[anchor=south east, font=\sffamily\fontsize{8.5}{9}\selectfont, text=bkInk, inner sep=0pt] at (-0.4,#3*#4+2.2) {（#7）};
\end{scope}}
% ぼう1本  \bkqpbar{(x,0)}{はば}{高さ mm}{色}   (x,0) は ぼうの下のまん中
\newcommand\bkqpbar[4]{\begin{scope}[shift={#1}]\colorlet{bkqp@b}{#4}%
  \filldraw[fill=bkqp@b, draw=bkqp@b!60!black, line width=0.5pt] ({-(#2)/2},0) rectangle ({(#2)/2},{#3});\end{scope}}
% ぼうグラフ  \bkqpgraph{(x,y)}{はば}{目もりの数}{1目もりの高さ}{数字を書く間}{1目もりの数}{単位}{ぼうの間}{数/名前,…}
\newcommand\bkqpbarcolor{bkBlue!55}
\newcommand\bkqpgraph[9]{\begin{scope}[shift={#1}]
  \bkqpaxes{(0,0)}{#2}{#3}{#4}{#5}{#6}{#7}
  \foreach \v/\t [count=\i] in {#9} {
    \pgfmathsetmacro\bkqp@h{\v/#6*#4}
    \bkqpbar{({(\i-0.5)*#8},0)}{#8*0.56}{\bkqp@h}{\bkqpbarcolor}
    \node[anchor=base, font=\sffamily\bfseries\fontsize{8.5}{10}\selectfont, text=bkInk, inner sep=0pt] at ({(\i-0.5)*#8},-4.6) {\t};}
\end{scope}}
% グラフの下の名前だけ  \bkqpnames{ぼうの間}{名前,…}
\newcommand\bkqpnames[2]{\foreach \t [count=\i] in {#2} {
    \node[anchor=base, font=\sffamily\bfseries\fontsize{8.5}{10}\selectfont, text=bkInk, inner sep=0pt] at ({(\i-0.5)*#1},-4.6) {\t};}}
% グラフの題名  \bkqptitle{(x,y)}{題名}
\newcommand\bkqptitle[2]{\node[font=\sffamily\bfseries\fontsize{9.5}{11}\selectfont, text=bkInk, inner sep=0pt] at #1 {#2};}
% 【第11回】正の字の線だけ（tikz の中で）  \bkqpstrokes{(左下)}{1字の大きさ}{本数 1〜5}
\newcommand\bkqpstrokes[3]{\bkqp@strokes{#1}{#2}{#3}}
% 【第11回】正の字の書きじゅん（番号つき）  \bkqpstrokeorder{(左下)}{1字の大きさ}
\newcommand\bkqpstrokeorder[2]{\begin{scope}[shift={#1}]
  \bkqp@strokes{(0,0)}{#2}{5}
  \foreach \n/\px/\py in {1/{0.5*#2}/{#2+3.4}, 2/{0.5*#2-3}/{0.84*#2}, 3/{0.86*#2+3.2}/{0.52*#2}, 4/{0.17*#2-3.2}/{0.3*#2}, 5/{0.5*#2}/{-3.4}} {
    \node[circle, fill=bkTheme, text=white, font=\sffamily\bfseries\fontsize{8}{8}\selectfont, inner sep=0pt, minimum size=4.2mm] at (\px,\py) {\n};}
\end{scope}}
% 【第11回】目もりを 1つずつ かぞえて 上がる 矢印  \bkqpclimb{(x,0)}{1目もりの高さ}{いくつ}{1目もりの数}{色}
\newcommand\bkqpclimb[5]{\begin{scope}[shift={#1}]
  \foreach \k in {1,...,#3} {
    \draw[#5, line width=0.8pt, -{Stealth[length=2.6pt]}] (0,{(\k-1)*#2+0.3}) .. controls (2.2,{(\k-1)*#2+0.3}) and (2.2,{\k*#2-0.3}) .. (0.3,{\k*#2-0.3});
    \pgfmathtruncatemacro\bkqp@cv{\k*#4}
    \node[anchor=west, font=\sffamily\bfseries\fontsize{7.5}{8}\selectfont, text=#5, inner sep=0pt] at (2.6,{(\k-0.5)*#2+0.6}) {\bkqp@cv};}
\end{scope}}
\makeatother

\bkchapter{1}{17}{\R{表}{ひょう}と ぼうグラフ}{\bklevel{1}}{ひょうと ぼうグラフ}
\bkmichishirube{1-5,1-11,1-13}{1-18,3-3,3-13}

\begin{goal}
\begin{itemize}[label={\color{bkTheme}●}, leftmargin=1.4\zw, itemsep=0pt, topsep=0pt]
\item しらべて かぞえ、 \R{表}{ひょう}に まとめられる。
\item ぼうグラフが かけて、 よめる。
\item グラフの 1\R{目}{め}もりを たしかめられる。
\end{itemize}
\end{goal}

\begin{hatena}
\begin{bkfigure}
\begin{tikzpicture}[x=1mm, y=1mm]
  % こくばん
  \fill[bkStem!80, rounded corners=2mm] (-2,-2) rectangle (124,62);
  \fill[bkGreen!55!black, rounded corners=1mm] (0,0) rectangle (122,60);
  \fill[bkStem!60] (8,-4.5) rectangle (114,-2);
  \fill[white] (96,-3.2) rectangle (103,-2.2); \fill[bkYellow!70] (105,-3.2) rectangle (111,-2.2);
  \node[anchor=west, text=white, font=\sffamily\bfseries\small] at (5,54.5) {すきな きゅう\R{食}{しょく}を 1つ かいてね};
  \foreach \x/\y/\a/\k in {16.8/45.8/-3/3,39.2/46.1/-7/4,60.9/46.1/5/2,85.9/46.2/7/1,106.0/46.2/7/1,15.0/36.5/-3/1,39.8/36.4/-7/1,60.1/36.2/-7/2,84.1/36.3/-3/3,106.2/36.4/4/3,13.3/27.1/-7/1,37.8/27.2/-7/1,62.7/27.0/-5/4,83.8/27.1/4/2,107.6/27.3/3/4,16.7/17.7/-7/3,38.2/17.6/-6/1,61.4/17.5/-5/2,82.9/17.3/-5/1,107.2/17.3/-6/2,15.5/8.3/-5/2,39.0/7.7/-6/4,59.6/8.1/5/3,83.0/8.0/3/2} {
    \bkqpcard{(\x,\y)}{\a}{\k}}
  \node[text=white, font=\sffamily\bfseries\fontsize{20}{20}\selectfont, rotate=8] at (107,8.5) {？};
\end{tikzpicture}
\end{bkfigure}
\noindent クラスの みんなが、 すきな きゅう\R{食}{しょく}を カードに かきました。\\
いちばん \R{人気}{にん|き}の きゅう\R{食}{しょく}は どれ？ \R{一目}{ひと|め}で わかるかな？
\end{hatena}

\clearpage
%% ---------------- p.173 やってみよう ----------------
\begin{yattemiyou}[カードを かぞえて、 \R{表}{ひょう}に しよう]
\begin{enumerate}[label={\bkcircled[bkGreen]{\arabic*}}, leftmargin=2.2\zw, itemsep=2pt, topsep=0pt]
\item \bkhatenaref の カードを 1まい かぞえたら、 その カードに ／を つけよう。
\item かぞえた きゅう\R{食}{しょく}の よこに、 \R{線}{せん}を 1\R{本}{ぽん} かこう。 \R{線}{せん}は 「\R{正}{せい}」の \R{字}{じ}の じゅんに かきます。
\end{enumerate}
\begin{bkfigure}
\begin{tikzpicture}[x=1mm, y=1mm]
  \foreach \n in {1,...,5} {
    \node[inner sep=0pt] at ({(\n-1)*22},0) {\bkqpsei[7]{\n}};
    \node[font=\sffamily\bfseries\small, text=bkGreen!70!black] at ({(\n-1)*22},-8) {\n\R{人}{にん}};}
  \foreach \n in {1,...,4} {\draw[bk arrow, bkGreen!60, line width=1pt] ({(\n-1)*22+6},0) -- ({\n*22-6},0);}
\end{tikzpicture}
\end{bkfigure}
\begin{enumerate}[label={\bkcircled[bkGreen]{\arabic*}}, leftmargin=2.2\zw, itemsep=2pt, topsep=0pt, start=3]
\item ぜんぶ かぞえたら、 \R{正}{せい}の \R{字}{じ}を \R{数}{かず}に して かこう。
\end{enumerate}
\par\vspace{3pt}
\begin{bkcenter}
\begin{tblr}{colspec={Q[l,m,wd=36mm] Q[l,m,wd=50mm] Q[c,m,wd=26mm]}, hlines={0.6pt,bkGreen!45}, vlines={0.6pt,bkGreen!45},
  row{1}={bg=bkGreen!15, font=\sffamily\bfseries\small}, rowsep=5pt, colsep=5pt, cells={valign=m}, row{2-6}={bg=white}}
\SetCell{c} きゅう\R{食}{しょく} & \SetCell{c} \R{正}{せい}の \R{字}{じ} & \R{人数}{にん|ずう} \\
\tikz[x=1mm,y=1mm,baseline=-1mm]{\bkqpfood{(0,0)}{1}{1}}\enspace カレー & & \bkkakikomi\,\R{人}{にん} \\
\tikz[x=1mm,y=1mm,baseline=-1mm]{\bkqpfood{(0,0)}{1}{2}}\enspace あげパン & & \bkkakikomi\,\R{人}{にん} \\
\tikz[x=1mm,y=1mm,baseline=-1mm]{\bkqpfood{(0,0)}{1}{3}}\enspace ラーメン & & \bkkakikomi\,\R{人}{にん} \\
\tikz[x=1mm,y=1mm,baseline=-1mm]{\bkqpfood{(0,0)}{1}{4}}\enspace やきそば & & \bkkakikomi\,\R{人}{にん} \\
\SetCell{c} ぜんぶ & & \bkkakikomi\,\R{人}{にん} \\
\end{tblr}
\end{bkcenter}
\noindent いちばん \R{人気}{にん|き}の きゅう\R{食}{しょく}は？\enspace\bkkakikomi[28mm]
\end{yattemiyou}

\bkfukidashi{mina}{／を つけると、 2かい かぞえないね。 まちがえても、 かぞえなおせば だいじょうぶ！}

\clearpage
%% ---------------- 数学の道具 ----------------
\bksection{doguu}{\R{数学}{すう|がく}の \R{道具}{どう|ぐ}}

\bksubsection{1}{かぞえて、 \R{表}{ひょう}に まとめる}

\label{bk:1-17:hyou}\bkhatenaref の カードは、 ばらばらに はって あります。
この ままでは、 どれが いくつ あるか、 すぐには わかりません。
そこで、 やってみようでは \R{正}{せい}の \R{字}{じ}\bkidx{せいのじ}{\R{正}{せい}の\R{字}{じ}}を つかって かぞえました。

\par\vspace{2pt}
\noindent{\sffamily\bfseries\color{bkTheme}\R{正}{せい}の \R{字}{じ}の ひみつ}

\noindent\begin{minipage}[c]{0.36\linewidth}\centering
\begin{tikzpicture}[x=1mm, y=1mm]
  \bkqpstrokeorder{(0,0)}{18}
\end{tikzpicture}
\end{minipage}\hfill
\begin{minipage}[c]{0.62\linewidth}
「\R{正}{せい}」の \R{字}{じ}は、 ちょうど 5\R{本}{ほん}の \R{線}{せん}で できて います。
\R{線}{せん}を 1\R{本}{ぽん} かく たびに、 1\R{人}{にん} かぞえます。
だから、 \R{正}{せい}が 1つ できたら 5\R{人}{にん}です。
\end{minipage}

\par\vspace{4pt}
\noindent\begin{minipage}[c]{0.36\linewidth}\centering
\begin{tikzpicture}[x=1mm, y=1mm]
  \bkqpstrokes{(0,0)}{10}{5}
  \bkqpstrokes{(15,0)}{10}{3}
  \bkbrace{(0,-2)}{(10,-2)}{5}{-1}
  \bkbrace{(15,-2)}{(25,-2)}{3}{-1}
\end{tikzpicture}
\end{minipage}\hfill
\begin{minipage}[c]{0.62\linewidth}
カレーは 「\R{正}{せい}」が 1つと、 3\R{本}{ぼん}。\\
\R{正}{せい}で 5、 のこりの 3\R{本}{ぼん}で 3。\\
5 ＋ 3 ＝ 8 で、 8\R{人}{にん}です。
\end{minipage}

\par\vspace{4pt}
\noindent 5の まとまりが できるので、 1\R{本}{ぽん}ずつ かぞえなおさなくても すみます。
10の まとまり（1-2）と \R{同}{おな}じ くふうです。

\noindent ぜんぶ かぞえたら、 \R{表}{ひょう}\bkidx{ひょう}{\R{表}{ひょう}}（1-13）に まとめます。

\begin{bkcenter}
\begin{bktable}{colspec={Q[c,wd=21mm] Q[c,wd=18mm] Q[c,wd=18mm] Q[c,wd=18mm] Q[c,wd=18mm] Q[c,wd=14mm]}, colsep=3pt, cell{1-3}{1}={bg=bkTheme!12, font=\sffamily\bfseries\small}, cell{1-3}{6}={bg=bkTheme!5}}
きゅう\R{食}{しょく} & カレー & あげパン & ラーメン & やきそば & ぜんぶ \\
\R{正}{せい}の \R{字}{じ} & \bkqpsei{8} & \bkqpsei{7} & \bkqpsei{5} & \bkqpsei{4} & \\
\R{人数}{にん|ずう} & 8\R{人}{にん} & 7\R{人}{にん} & 5\R{人}{にん} & 4\R{人}{にん} & 24\R{人}{にん} \\
\end{bktable}
\end{bkcenter}

\noindent 「ぜんぶ」の \R{人数}{にん|ずう}は、 4つの \R{人数}{にん|ずう}を たして もとめます。
\begin{enumerate}[label={\bkcircled[bkTheme]{\arabic*}}, leftmargin=2.2\zw, itemsep=0pt, topsep=2pt]
\item カレーと あげパン\hspace{1\zw}8 ＋ 7 ＝ 15
\item ラーメンを たす\hspace{1\zw}15 ＋ 5 ＝ 20
\item やきそばを たす\hspace{1\zw}20 ＋ 4 ＝ 24
\end{enumerate}

\noindent\begin{minipage}[c]{0.62\linewidth}
\bkhatenaref の カードを 1まいずつ かぞえると、 24まい あります。
\R{表}{ひょう}の 「ぜんぶ」の 24\R{人}{にん}と ぴったり \R{同}{おな}じです。
だから、 かぞえわすれた カードも、 2かい かぞえた カードも ない ことが わかります。
\end{minipage}\hfill
\begin{minipage}[c]{0.34\linewidth}\centering
\begin{tikzpicture}[x=1mm, y=1mm]
  \foreach \i in {0,...,23} {\pgfmathtruncatemacro\cx{mod(\i,6)}\pgfmathtruncatemacro\cy{\i/6}
    \filldraw[fill=white, draw=bkGray!80, line width=0.5pt, rounded corners=0.4mm] ({\cx*6.4},{-\cy*5}) rectangle ++(5.4,3.8);
    \draw[bkTheme, line width=0.7pt] ({\cx*6.4+1.4},{-\cy*5+0.8}) -- ++(2.6,2.2);}
  \node[font=\sffamily\bfseries\small, text=bkTheme] at (18.6,-20.5) {24まい};
\end{tikzpicture}
\end{minipage}

\bksubsection{2}{ぼうグラフに する}

\noindent \R{表}{ひょう}を \R{見}{み}れば、 \R{人数}{にん|ずう}は わかります。
でも、 8、 7、 5、 4 と \R{数}{かず}を 1つずつ \R{見}{み}くらべないと、 どれが \R{多}{おお}いかは わかりません。
そこで、 \R{人数}{にん|ずう}を ぼうの \R{高}{たか}さに して ならべて みましょう。

\noindent\begin{minipage}[c]{0.6\linewidth}\centering\label{bk:1-17:bou}%
\begin{tikzpicture}[x=1mm, y=1mm]
  \bkqptitle{(36,43)}{すきな きゅう\R{食}{しょく}}
  \bkqpaxes{(0,0)}{68}{10}{3.6}{1}{1}{\R{人}{にん}}
  \foreach \v/\k [count=\i] in {8/1,7/2,5/3,4/4} {
    \bkqpbar{({(\i-0.5)*17},0)}{10}{\v*3.6}{bkBlue!55}
    \bkqpfood{({(\i-0.5)*17},-4.2)}{0.85}{\k}
    \node[anchor=north, font=\sffamily\bfseries\fontsize{8.5}{10}\selectfont, text=bkInk, inner sep=0pt] at ({(\i-0.5)*17},-7.4) {\bkqpname{\k}};}
  \foreach \t/\x/\y in {1/-6.5/-6.5,2/-10/20,3/16/31.5,4/18/43} {
    \node[circle, fill=bkTheme, text=white, font=\sffamily\bfseries\fontsize{8.5}{9}\selectfont, inner sep=0pt, minimum size=4.4mm] at (\x,\y) {\t};}
\end{tikzpicture}
\end{minipage}\hfill
\begin{minipage}[c]{0.38\linewidth}
\begin{enumerate}[label={\bkcircled[bkTheme]{\arabic*}}, leftmargin=2\zw, itemsep=3pt, topsep=0pt]
\item よこに \R{名前}{な|まえ}
\item たてに \R{目}{め}もりと \R{数}{かず}
\item ぼうを ぬる
\item \R{上}{うえ}に グラフの \R{名前}{な|まえ}
\end{enumerate}
\end{minipage}

\par\vspace{6pt}
\noindent{\sffamily\bfseries\color{bkTheme}ぼうの \R{高}{たか}さの きめかた}

\noindent\begin{minipage}[c]{0.36\linewidth}\centering
\begin{tikzpicture}[x=1mm, y=1mm]
  \bkqpaxes{(0,0)}{24}{10}{4}{1}{1}{\R{人}{にん}}
  \bkqpbar{(9,0)}{10}{8*4}{bkBlue!55}
  \bkqpclimb{(14.5,0)}{4}{8}{1}{bkTheme}
  \draw[bkTheme, line width=0.8pt, dash pattern=on 2pt off 1.5pt] (4,32) -- (22,32);
  \node[anchor=north, font=\sffamily\bfseries\fontsize{8.5}{10}\selectfont, text=bkInk, inner sep=0pt] at (9,-1.6) {カレー};
\end{tikzpicture}
\end{minipage}\hfill
\begin{minipage}[c]{0.62\linewidth}
カレーは 8\R{人}{にん}。 この グラフの 1\R{目}{め}もりは 1\R{人}{にん}です。
\begin{enumerate}[label={\bkcircled[bkTheme]{\arabic*}}, leftmargin=2\zw, itemsep=0pt, topsep=2pt]
\item 0の \R{線}{せん}から はじめる。
\item \R{目}{め}もりを 1、 2、 3、 … と 1つずつ かぞえて \R{上}{あ}がる。
\item 8の \R{線}{せん}まで \R{来}{き}たら、 そこまで ぬる。
\end{enumerate}
ものさしで 0から はかるのと \R{同}{おな}じです（1-3、 1-8）。
\end{minipage}

\par\vspace{6pt}
\noindent \R{人数}{にん|ずう}を ぼうの \R{高}{たか}さで あらわした グラフを、 ぼうグラフ\bkidx{ぼうぐらふ}{ぼうグラフ}と いいます。\\
\bkhatenaref の \R{答}{こた}えは カレー。 ぼうが いちばん \R{高}{たか}いので、 \R{一目}{ひと|め}で わかります。

\bksubsection{3}{\R{表}{ひょう}と ぼうグラフの よい ところ}

\noindent \R{同}{おな}じ しらべた ことを、 \R{表}{ひょう}と ぼうグラフに しました。 2つの しつもんに \R{答}{こた}えて みよう。

\par\vspace{3pt}
\noindent\begin{minipage}[c]{0.5\linewidth}\centering
\begin{bktable}{colspec={Q[c,wd=19mm] Q[c,wd=11mm]}, colsep=2pt, rowsep=1pt, row{1}={bg=bkTheme!12, font=\sffamily\bfseries\small}, cells={font=\small}}
きゅう\R{食}{しょく} & \R{人数}{にん|ずう} \\
カレー & 8 \\
あげパン & 7 \\
ラーメン & 5 \\
やきそば & 4 \\
\end{bktable}
\end{minipage}\hfill
\begin{minipage}[c]{0.48\linewidth}\centering
\begin{tikzpicture}[x=1mm, y=1mm]
  \bkqpgraph{(0,0)}{52}{10}{2.8}{2}{1}{\R{人}{にん}}{13}{8/カレー,7/あげパン,5/ラーメン,4/やきそば}
\end{tikzpicture}
\end{minipage}

\par\vspace{6pt}
\noindent\begin{tblr}{colspec={Q[l,m,wd=50mm] Q[l,m,wd=40mm] Q[l,m,wd=40mm]}, hlines={0.6pt,bkGray!50}, vlines={0.6pt,bkGray!50}, rowsep=4pt, colsep=4pt, row{1}={bg=bkTheme!10, font=\sffamily\bfseries\small}, cells={font=\small}}
しつもん & \R{表}{ひょう}で さがすと & グラフで さがすと \\
あげパンは \R{何人}{なん|にん}？ & 7と すぐ わかる\enspace\bkmaru & \R{目}{め}もりを かぞえる \\
2ばんめに \R{多}{おお}いのは？ & 4つの \R{数}{かず}を くらべる & 2ばんめに \R{高}{たか}い ぼう\enspace\bkmaru \\
\end{tblr}

\par\vspace{6pt}
\noindent \R{表}{ひょう}は、 1つ 1つの \R{数}{かず}が \R{正}{ただ}しく わかります。
ぼうグラフは、 どれが \R{多}{おお}いか、 \R{一目}{ひと|め}で くらべられます。
つかい\R{方}{かた}に あわせて、 えらびましょう。

\begin{taisetsu}
\R{表}{ひょう}は、 \R{数}{かず}が \R{正}{ただ}しく わかる。\\
ぼうグラフは、 \R{一目}{ひと|め}で くらべられる。
\end{taisetsu}

\bksubsection{4}{ぼうグラフを よむ}

\noindent\begin{minipage}[c]{0.6\linewidth}
\label{bk:1-17:yomu}いちばん \R{多}{おお}いのは、 いちばん \R{高}{たか}い カレーで 8\R{人}{にん}。
いちばん \R{少}{すく}ないのは、 いちばん ひくい やきそばで 4\R{人}{にん}。

\vspace{3pt}
カレーと やきそばの ちがいを \R{見}{み}よう。
やきそばの ぼうの \R{上}{うえ}から よこに \R{線}{せん}を ひくと、 カレーの ぼうが はみ\R{出}{だ}します。
はみ\R{出}{だ}した ところが ちがいです。
\end{minipage}\hfill
\begin{minipage}[c]{0.37\linewidth}\centering
\begin{tikzpicture}[x=1mm, y=1mm]
  \bkqpaxes{(0,0)}{38}{10}{3.4}{1}{1}{\R{人}{にん}}
  \bkqpbar{(10,0)}{10}{8*3.4}{bkTheme!75}
  \bkqpbar{(28,0)}{10}{4*3.4}{bkBlue!55}
  \draw[bkTheme, line width=0.8pt, dash pattern=on 2pt off 1.5pt] (5,4*3.4) -- (36,4*3.4);
  \draw[bkTheme, line width=0.8pt, dash pattern=on 2pt off 1.5pt] (15,8*3.4) -- (36,8*3.4);
  \draw[bkTheme, line width=1pt, decorate, decoration={brace, amplitude=1.5mm}] (35.5,8*3.4) -- (35.5,4*3.4)
    node[midway, right=2.3mm, font=\sffamily\bfseries\small, text=bkTheme, inner sep=0pt] {4};
  \bkqpnames{18}{カレー,やきそば}
\end{tikzpicture}
\end{minipage}

\par\vspace{4pt}
\noindent ちがいは、 2つの やり\R{方}{かた}で もとめられます。
\begin{enumerate}[label={\bkcircled[bkTheme]{\arabic*}}, leftmargin=2.2\zw, itemsep=0pt, topsep=2pt]
\item グラフで：はみ\R{出}{だ}した ところの \R{目}{め}もりを かぞえる。 4\R{目}{め}もりで 4\R{人}{にん}。
\item \R{数}{かず}で：ひき\R{算}{ざん}（1-5）で、 8 － 4 ＝ 4。 4\R{人}{にん}。
\end{enumerate}
どちらで やっても、 \R{同}{おな}じ 4\R{人}{にん}に なります。

\bksubsection{5}{1\R{目}{め}もりの \R{大}{おお}きさ}

\label{bk:1-17:memori}\R{下}{した}の ぼうは、 どれも 4\R{目}{め}もり。
でも、 1\R{目}{め}もり\bkidx{ひとめもり}{1\R{目}{め}もり}の \R{大}{おお}きさが ちがうと、 \R{人数}{にん|ずう}も ちがいます。

\begin{bkfigure}
\begin{tikzpicture}[x=1mm, y=1mm]
  \foreach \u [count=\i] in {1,2,5,10} {
    \begin{scope}[shift={({(\i-1)*33+8},0)}]
      \bkqpaxes{(0,0)}{16}{5}{3.8}{1}{\u}{\R{人}{にん}}
      \bkqpbar{(8,0)}{9}{4*3.8}{bkBlue!55}
      \pgfmathtruncatemacro\n{4*\u}
      \node[anchor=base, font=\sffamily\bfseries\small, text=bkBlue!70!black] at (8,17.5) {\n\R{人}{にん}};
      \node[anchor=north, font=\sffamily\bfseries\fontsize{9}{10}\selectfont, text=bkTheme, inner sep=0pt] at (6,-2.4) {1\R{目}{め}もり \u\R{人}{にん}};
    \end{scope}}
\end{tikzpicture}
\end{bkfigure}

\noindent 1\R{目}{め}もり 2\R{人}{にん}の グラフでは、 \R{目}{め}もりを 1つ \R{上}{あ}がる たびに 2\R{人}{にん}ずつ ふえます。
2、 4、 6、 8 と かぞえて、 4\R{目}{め}もりで 8\R{人}{にん}です。
2が 4つ なので、 かけ\R{算}{ざん}（1-12）で 2 × 4 ＝ 8 とも かけます。

\par\vspace{3pt}
\noindent\begin{tblr}{colspec={Q[c,m,wd=26mm] Q[l,m,wd=62mm] Q[c,m,wd=32mm]}, hlines={0.6pt,bkGray!50}, vlines={0.6pt,bkGray!50}, rowsep=3pt, colsep=4pt, row{1}={bg=bkTheme!10, font=\sffamily\bfseries\small}, cells={font=\small}}
1\R{目}{め}もり & 4\R{目}{め}もりまで かぞえる & 4\R{目}{め}もりは \\
1\R{人}{にん} & 1、 2、 3、 4 & 1 × 4 ＝ 4 \\
2\R{人}{にん} & 2、 4、 6、 8 & 2 × 4 ＝ 8 \\
5\R{人}{にん} & 5、 10、 15、 20 & 5 × 4 ＝ 20 \\
10\R{人}{にん} & 10、 20、 30、 40 & 10が 4つで 40 \\
\end{tblr}

\noindent{\sffamily\bfseries\color{bkTheme}1\R{目}{め}もりの \R{見}{み}つけかた}

\noindent グラフには、 \R{数}{かず}が すべての \R{目}{め}もりに かいて あるとは かぎりません。
\R{数}{かず}が かいて ある \R{目}{め}もりを 2つ \R{見}{み}つけて、 その \R{間}{あいだ}を しらべます。

\par\vspace{3pt}
\noindent\begin{minipage}[c]{0.36\linewidth}\centering
\begin{tikzpicture}[x=1mm, y=1mm]
  \bkqpaxes{(0,0)}{20}{10}{3.6}{5}{2}{\R{人}{にん}}
  \bkqpclimb{(21,0)}{3.6}{5}{2}{bkTheme}
\end{tikzpicture}
\end{minipage}\hfill
\begin{minipage}[c]{0.62\linewidth}
\begin{enumerate}[label={\bkcircled[bkTheme]{\arabic*}}, leftmargin=2\zw, itemsep=0pt, topsep=0pt]
\item \R{数}{かず}が かいて あるのは 0、 10、 20。
\item 0から 10までに、 \R{目}{め}もりは 5つ。
\item 10を 5つに \R{同}{おな}じずつ \R{分}{わ}けると 2。\\
2 ＋ 2 ＋ 2 ＋ 2 ＋ 2 ＝ 10 だからです。\\
わり\R{算}{ざん}（1-14）で 10 ÷ 5 ＝ 2 とも かけます。
\item だから 1\R{目}{め}もりは 2\R{人}{にん}。
\end{enumerate}
\end{minipage}

\begin{taisetsu}
グラフを よむ ときは、 まず 1\R{目}{め}もりの \R{大}{おお}きさを たしかめる。
\end{taisetsu}

\bksubsection{6}{1\R{目}{め}もりの きめかた}

\label{bk:1-17:kimeru}ペットボトルの ふたを あつめました。 1\R{組}{くみ} 40こ、 2\R{組}{くみ} 20こ、 3\R{組}{くみ} 30こ、 4\R{組}{くみ} 10こ。
10\R{目}{め}もりの グラフに かきます。 1\R{目}{め}もりを いくつに すると よいでしょう。

\par\vspace{3pt}
\noindent 1\R{目}{め}もりを きめると、 10\R{目}{め}もりで いくつまで かけるかが きまります。

\begin{bkcenter}
\begin{tblr}{colspec={Q[c,m,wd=24mm] Q[c,m,wd=44mm] Q[c,m,wd=40mm]}, hlines={0.6pt,bkGray!50}, vlines={0.6pt,bkGray!50}, rowsep=2pt, colsep=4pt, row{1}={bg=bkTheme!10, font=\sffamily\bfseries\small}, cells={font=\small}}
1\R{目}{め}もり & 10\R{目}{め}もりで かける \R{数}{かず} & 40こは \R{入}{はい}る？ \\
2こ & 2が 10こで 20こ まで & \R{入}{はい}らない \\
5こ & 5が 10こで 50こ まで & \R{入}{はい}る \\
10こ & 10が 10こで 100こ まで & \R{入}{はい}る \\
\end{tblr}
\end{bkcenter}

\begin{bkfigure}
\begin{tikzpicture}[x=1mm, y=1mm]
  % 1目もり 2こ：はみ出す
  \begin{scope}[shift={(8,0)}]
    \bkqpaxes{(0,0)}{30}{10}{3}{5}{2}{こ}
    \foreach \cx in {3.75,18.75} {
      \fill[bkBlue!55] (\cx-2.5,0) rectangle (\cx+2.5,34);
      \draw[bkBlue!55!black, line width=0.5pt] (\cx-2.5,34) -- (\cx-2.5,0) -- (\cx+2.5,0) -- (\cx+2.5,34);
      \draw[white, line width=1.6pt] (\cx-3.3,31.6) -- (\cx-1.2,32.8) -- (\cx+1.2,31.6) -- (\cx+3.3,32.8);}
    \node[font=\sffamily\bfseries\fontsize{7.5}{8}\selectfont, text=bkRed, inner sep=0pt] at (12,38) {はみ\R{出}{だ}す！};
    \bkqpbar{(11.25,0)}{5}{30}{bkBlue!55}
    \bkqpbar{(26.25,0)}{5}{15}{bkBlue!55}
    \node[anchor=north, font=\sffamily\bfseries\fontsize{8.5}{10}\selectfont, text=bkTheme, inner sep=0pt] at (15,-2) {1\R{目}{め}もり 2こ};
    \node at (15,-9.5) {\bkbatsu};
  \end{scope}
  % 1目もり 5こ：ちょうど よい
  \begin{scope}[shift={(51,0)}]
    \bkqpaxes{(0,0)}{30}{10}{3}{2}{5}{こ}
    \foreach \v [count=\i] in {40,20,30,10} {\bkqpbar{({(\i-0.5)*7.5},0)}{5}{\v/5*3}{bkBlue!55}}
    \node[anchor=north, font=\sffamily\bfseries\fontsize{8.5}{10}\selectfont, text=bkTheme, inner sep=0pt] at (15,-2) {1\R{目}{め}もり 5こ};
    \node at (15,-9.5) {\bkmaru};
  \end{scope}
  % 1目もり 10こ：ひくい
  \begin{scope}[shift={(94,0)}]
    \bkqpaxes{(0,0)}{30}{10}{3}{2}{10}{こ}
    \foreach \v [count=\i] in {40,20,30,10} {\bkqpbar{({(\i-0.5)*7.5},0)}{5}{\v/10*3}{bkBlue!55}}
    \node[anchor=north, font=\sffamily\bfseries\fontsize{8.5}{10}\selectfont, text=bkTheme, inner sep=0pt] at (15,-2) {1\R{目}{め}もり 10こ};
    \node[font=\sffamily\bfseries\fontsize{8}{9}\selectfont, text=bkOrange!80!black, align=center, inner sep=0pt] at (15,-9.5) {ぼうが ひくくて\\くらべにくい};
  \end{scope}
\end{tikzpicture}
\end{bkfigure}

\noindent 1\R{目}{め}もり 2こ では、 40こ や 30こ の ぼうが グラフから はみ\R{出}{だ}して しまいます。
1\R{目}{め}もり 10こ でも かけますが、 ぼうが ひくくて、 ちがいが \R{見}{み}えにくく なります。
\R{入}{はい}る ものの \R{中}{なか}で いちばん \R{小}{ちい}さい 1\R{目}{め}もり（5こ）に すると、 ぼうが \R{高}{たか}く なって くらべやすく なります。

\par\vspace{3pt}
\noindent 1\R{目}{め}もり 5こ なら、 40こは 5が 8つ（5、 10、 15、 … 40）なので 8\R{目}{め}もりです。

\begin{taisetsu}
1\R{目}{め}もりは、 いちばん \R{大}{おお}きい \R{数}{かず}が \R{入}{はい}る ものの \R{中}{なか}から、 いちばん \R{小}{ちい}さい ものを えらぶ。
\end{taisetsu}

%% ---------------- 例題 ----------------

\begin{reidai}[\R{正}{せい}の \R{字}{じ}を \R{数}{かず}に する]
すきな どうぶつを \R{正}{せい}の \R{字}{じ}で かぞえました。 \R{何人}{なん|にん}？\par
\noindent\hspace*{1\zw}イヌ\enspace\bkqpsei{10}\hspace{4\zw}ネコ\enspace\bkqpsei{8}

\kangaekata \R{正}{せい} 1つで 5。\\
イヌは \R{正}{せい}が 2つ。 5 ＋ 5 ＝ 10。\\
ネコは \R{正}{せい}が 1つと 3\R{本}{ぼん}。 5 ＋ 3 ＝ 8。

\kotae イヌ 10\R{人}{にん}、 ネコ 8\R{人}{にん}

\noindent はんたいに、 7\R{人}{にん}を \R{正}{せい}の \R{字}{じ}で かくと？\\
7 ＝ 5 ＋ 2。 \R{正}{せい}を 1つ かいて、 のこり 2\R{本}{ほん}。\enspace → \bkqpsei{7}
\end{reidai}

\begin{jibun}
\R{正}{せい}の \R{字}{じ}を \R{数}{かず}に しよう。\par
\noindent\hspace*{1\zw}イチゴ\enspace\bkqpsei{12}\enspace → \bkkakikomi\,\R{人}{にん}\hspace{2\zw}ミカン\enspace\bkqpsei{9}\enspace → \bkkakikomi\,\R{人}{にん}
\par\noindent 6\R{人}{にん}を \R{正}{せい}の \R{字}{じ}で かこう。\enspace\bkkakikomi[20mm]
\bkjans{イチゴ 12\R{人}{にん}（5 ＋ 5 ＋ 2）、 ミカン 9\R{人}{にん}（5 ＋ 4）。 6\R{人}{にん}は \R{正}{せい}と 1\R{本}{ぽん}（6 ＝ 5 ＋ 1）}
\end{jibun}

\begin{reidai}[ぼうグラフを よむ]
サッカーと てつぼうの ちがいは \R{何人}{なん|にん}？\par
\noindent\begin{minipage}[c]{0.56\linewidth}\centering
\begin{tikzpicture}[x=1mm, y=1mm]
  \bkqptitle{(32,33.5)}{すきな あそび}
  \bkqpgraph{(0,0)}{64}{10}{2.8}{2}{1}{\R{人}{にん}}{16}{9/サッカー,7/おにごっこ,6/なわとび,4/てつぼう}
\end{tikzpicture}
\end{minipage}\hfill
\begin{minipage}[c]{0.42\linewidth}
\kangaekata 1\R{目}{め}もりは 1\R{人}{にん}（0と 2の \R{間}{あいだ}が 2\R{目}{め}もり）。
サッカー 9\R{人}{にん}、 てつぼう 4\R{人}{にん}。

\shiki 9 － 4 ＝ 5\enspace \kotae 5\R{人}{にん}
\end{minipage}
\end{reidai}

\begin{jibun}
（1）\R{上}{うえ}の グラフで、 おにごっこと てつぼうの ちがいは？
\bkshikiran
（2）サッカーと なわとびの ちがいは？
\bkshikiran
\bkjans{（1）7 － 4 ＝ 3\quad 3\R{人}{にん}\quad（2）9 － 6 ＝ 3\quad 3\R{人}{にん}}
\end{jibun}

\begin{reidai}[1\R{目}{め}もりを たしかめて よむ]
\R{図書室}{と|しょ|しつ}で かりた \R{本}{ほん}。 \R{水曜日}{すい|よう|び}は \R{何}{なん}さつ？\par\vspace{5mm}
\noindent\begin{minipage}[c]{0.46\linewidth}\centering
\begin{tikzpicture}[x=1mm, y=1mm]
  \bkqpgraph{(0,0)}{48}{10}{2.8}{5}{2}{さつ}{9.6}{12/{\R{月}{げつ}},8/{\R{火}{か}},16/{\R{水}{すい}},10/{\R{木}{もく}},6/{\R{金}{きん}}}
\end{tikzpicture}
\end{minipage}\hfill
\begin{minipage}[c]{0.52\linewidth}
\kangaekata\\
\R{数}{かず}は 0、 10、 20 しか かいて いない。 まず 1\R{目}{め}もりを しらべよう。
\end{minipage}\par\vspace{2pt}
\begin{enumerate}[label={\bkcircled[bkBlue]{\arabic*}}, leftmargin=2.2\zw, itemsep=0pt, topsep=0pt]
\item 0から 10までに \R{目}{め}もりが 5つ。 2が 5つで 10 なので、 1\R{目}{め}もりは 2さつ。
\item \R{水曜日}{すい|よう|び}の ぼうは 8\R{目}{め}もり。
\item 2、 4、 6、 … と 8つ かぞえて 16。（2 × 8 ＝ 16）
\end{enumerate}
\kotae 16さつ
\end{reidai}

\begin{jibun}
\R{上}{うえ}の グラフで、 \R{月曜日}{げつ|よう|び}は \R{何}{なん}さつ？\enspace\bkkakikomi\,さつ
\bkjans{12さつ（6\R{目}{め}もり。 2 × 6 ＝ 12）}
\end{jibun}

\begin{reidai}[\R{表}{ひょう}を ぼうグラフに する]
\noindent\begin{minipage}[c]{0.45\linewidth}
あきかんの \R{数}{かず}を、 10\R{目}{め}もりの グラフに かきます。
1\R{目}{め}もりは \R{何}{なん}こ？
\end{minipage}\hfill
\begin{minipage}[c]{0.53\linewidth}\raggedleft
\begin{bktable}{colspec={Q[c,wd=17mm] Q[c,wd=8.5mm] Q[c,wd=8.5mm] Q[c,wd=8.5mm] Q[c,wd=8.5mm]}, colsep=1.5pt, rowsep=2pt, row{1}={bg=bkBlue!12}, cell{1-2}{1}={bg=bkBlue!12}, cells={font=\small}}
\R{組}{くみ} & 1 & 2 & 3 & 4 \\
あきかん & 18 & 12 & 16 & 8 \\
\end{bktable}
\end{minipage}
\par\vspace{2pt}
\noindent\kangaekata いちばん \R{多}{おお}いのは 18こ。
\begin{enumerate}[label={\bkcircled[bkBlue]{\arabic*}}, leftmargin=2.2\zw, itemsep=0pt, topsep=0pt]
\item 1\R{目}{め}もり 1こ：10こ までしか かけない。 18こは \R{入}{はい}らない。
\item 1\R{目}{め}もり 2こ：2が 10こで 20。 20こ まで かける。 18こが \R{入}{はい}る。
\end{enumerate}
\noindent\begin{minipage}[c]{0.58\linewidth}
\begin{enumerate}[label={\bkcircled[bkBlue]{\arabic*}}, leftmargin=2.2\zw, itemsep=0pt, topsep=0pt, start=3]
\item 18は 2が 9つ（2 × 9 ＝ 18）なので、 1\R{組}{くみ}は 9\R{目}{め}もり。 12は 2が 6つで 6\R{目}{め}もり。
\end{enumerate}
\kotae 1\R{目}{め}もり 2こ
\end{minipage}\hfill
\begin{minipage}[c]{0.4\linewidth}\centering
\begin{tikzpicture}[x=1mm, y=1mm]
  \bkqpaxes{(0,0)}{42}{10}{2.6}{2}{2}{こ}
  \bkqpbar{(5.25,0)}{6}{9*2.6}{\bkqpbarcolor}\bkqpbar{(15.75,0)}{6}{6*2.6}{\bkqpbarcolor}
  \bkqpnames{10.5}{{1\R{組}{くみ}},{2\R{組}{くみ}},{3\R{組}{くみ}},{4\R{組}{くみ}}}
\end{tikzpicture}
\end{minipage}
\end{reidai}

\begin{jibun}
\R{上}{うえ}の \R{表}{ひょう}を \R{見}{み}て、 3\R{組}{くみ}と 4\R{組}{くみ}の ぼうを グラフに かこう。
\bkjans{3\R{組}{くみ}は 8\R{目}{め}もり（2 × 8 ＝ 16）、 4\R{組}{くみ}は 4\R{目}{め}もり（2 × 4 ＝ 8）の ぼう}
\end{jibun}

\begin{reidai}[1\R{目}{め}もりを えらぶ]
すきな くだものを しらべました。 10\R{目}{め}もりの グラフに かくとき、 1\R{目}{め}もりは 1\R{人}{にん}・2\R{人}{にん}・5\R{人}{にん}の どれに する？\par\vspace{2pt}
\begin{bkcenter}
\begin{bktable}{colspec={Q[c,wd=20mm] Q[c,wd=14mm] Q[c,wd=14mm] Q[c,wd=14mm] Q[c,wd=14mm]}, colsep=1.5pt, rowsep=2pt, cell{1-2}{1}={bg=bkBlue!12, font=\sffamily\bfseries\small}, cells={font=\small}}
くだもの & りんご & みかん & ぶどう & もも \\
\R{人数}{にん|ずう} & 30 & 25 & 15 & 10 \\
\end{bktable}
\end{bkcenter}
\noindent\begin{minipage}[c]{0.58\linewidth}
\kangaekata いちばん \R{多}{おお}いのは 30\R{人}{にん}。
\begin{enumerate}[label={\bkcircled[bkBlue]{\arabic*}}, leftmargin=2\zw, itemsep=0pt, topsep=0pt]
\item 1\R{人}{にん}：10\R{人}{にん}まで。 \R{入}{はい}らない。
\item 2\R{人}{にん}：2が 10こで 20\R{人}{にん}まで。 \R{入}{はい}らない。
\item 5\R{人}{にん}：5が 10こで 50\R{人}{にん}まで。 \R{入}{はい}る。
\end{enumerate}
りんごは 5が 6つで 6\R{目}{め}もり（5 × 6 ＝ 30）。
\kotae 5\R{人}{にん}
\end{minipage}\hfill
\begin{minipage}[c]{0.4\linewidth}\centering
\begin{tikzpicture}[x=1mm, y=1mm]
  \bkqpgraph{(0,0)}{44}{10}{2.6}{2}{5}{\R{人}{にん}}{11}{30/りんご,25/みかん,15/ぶどう,10/もも}
\end{tikzpicture}
\end{minipage}
\end{reidai}

\begin{jibun}
あつめた どんぐりの \R{数}{かず}です。 10\R{目}{め}もりの グラフに かくとき、 1\R{目}{め}もりは 1こ・2こ・5この どれが よい？\par\vspace{2pt}
\begin{bkcenter}
\begin{bktable}{colspec={Q[c,wd=16mm] Q[c,wd=14mm] Q[c,wd=14mm] Q[c,wd=16mm] Q[c,wd=14mm]}, colsep=1.5pt, rowsep=2pt, cell{1-2}{1}={bg=bkTheme!12, font=\sffamily\bfseries\small}, cells={font=\small}}
\R{名前}{な|まえ} & ソラ & ミナ & デルタ & ケン \\
\R{数}{かず}（こ） & 16 & 10 & 6 & 14 \\
\end{bktable}
\end{bkcenter}
\noindent 1\R{目}{め}もり\enspace\bkkakikomi\,こ\hspace{2\zw}ソラの ぼうは\enspace\bkkakikomi\,\R{目}{め}もり
\bkjans{1\R{目}{め}もり 2こ（1こでは 10こ までで \R{入}{はい}らない。 2こなら 20こ まで \R{入}{はい}る。 5こでも \R{入}{はい}るが、 2この ほうが \R{小}{ちい}さい）。 ソラは 16 ＝ 2 × 8 で 8\R{目}{め}もり}
\end{jibun}

\par\vspace{6pt}
%% ---------------- p.178 まちがえやすいところ ----------------
\bksection{tsumazuki}{まちがえやすい ところ}

\begin{tsumazuki}{ぼうを とちゅうから かく}
どんぐりを ソラは 22こ、 ミナは 24こ ひろいました。
\begin{bkfigure}
\begin{tikzpicture}[x=1mm, y=1mm]
  \begin{scope}[shift={(10,0)}]
    \bkqpaxes{(0,0)}{34}{5}{5}{1}{0}{こ}
    \fill[bkRedBg] (-6,-1.8) rectangle (-1.1,1.8);
    \foreach \k/\v in {0/20,1/21,2/22,3/23,4/24,5/25} {\node[anchor=east, font=\sffamily\bfseries\fontsize{8.5}{9}\selectfont, text=bkInk, inner sep=0pt] at (-1.4,\k*5) {\v};}
    \bkqpbar{(9,0)}{10}{10}{bkBlue!55}\bkqpbar{(25,0)}{10}{20}{bkBlue!55}
    \bkqpnames{16}{ソラ,ミナ}
    \node[anchor=west] at (37,22) {\bkbatsu};
  \end{scope}
  \begin{scope}[shift={(72,0)}]
    \bkqpaxes{(0,0)}{34}{5}{5}{1}{5}{こ}
    \bkqpbar{(9,0)}{10}{22}{bkBlue!55}\bkqpbar{(25,0)}{10}{24}{bkBlue!55}
    \bkqpnames{16}{ソラ,ミナ}
    \node[anchor=west] at (37,22) {\bkmaru};
  \end{scope}
\end{tikzpicture}
\end{bkfigure}
\bkbatsu の グラフは、 20から かきはじめて います。 ソラの ぼうは 2\R{目}{め}もり、 ミナの ぼうは 4\R{目}{め}もり。
ミナが ソラの 2つ\R{分}{ぶん}も ひろった ように \R{見}{み}えます。
でも、 ほんとうは 24 － 22 ＝ 2で、 2こしか ちがいません。
ぼうの \R{高}{たか}さで \R{数}{かず}を くらべるので、 ぼうは かならず 0から かこう。
\end{tsumazuki}

\begin{tsumazuki}{1\R{目}{め}もりを たしかめない}
\noindent\begin{minipage}[c]{0.76\linewidth}
0から 10までに \R{目}{め}もりが 5つ。 1\R{目}{め}もりは 2\R{人}{にん}です。 ぼうは 6\R{目}{め}もり。\\
\bkbatsu\enspace 6\R{人}{にん}\hspace{3\zw}\bkmaru\enspace 12\R{人}{にん}\\
2、 4、 6、 8、 10、 12 と かぞえて 12\R{人}{にん}。 \R{目}{め}もりの \R{数}{かず}と \R{人数}{にん|ずう}は、 ちがう ことが あるよ。
\end{minipage}\hfill
\begin{minipage}[c]{0.22\linewidth}\centering
\begin{tikzpicture}[x=1mm, y=1mm]
  \bkqpaxes{(0,0)}{18}{8}{3}{5}{2}{\R{人}{にん}}
  \bkqpbar{(9,0)}{9}{18}{bkBlue!55}
\end{tikzpicture}
\end{minipage}
\end{tsumazuki}

\begin{tsumazuki}{ぜんぶの \R{数}{かず}を たしかめない}
カードは 24まい あるのに、 \R{表}{ひょう}の 「ぜんぶ」が 23\R{人}{にん}に なりました。\\
24 － 23 ＝ 1で、 どこかで 1まい かぞえわすれて います。
／の しるしが ついて いない カードを さがそう。
「ぜんぶ」と カードの \R{数}{かず}が \R{同}{おな}じに なるまで、 \R{表}{ひょう}は できあがりでは ありません。
\end{tsumazuki}

\begin{column}{\R{外国}{がい|こく}の かぞえ\R{方}{かた}}
\noindent\begin{minipage}[c]{0.79\linewidth}
えいごの \R{国}{くに}では、 たての \R{線}{せん}を 4\R{本}{ほん} かいて、 5\R{本}{ほん}\R{目}{め}で ななめに けします。
\R{正}{せい}の \R{字}{じ}と \R{同}{おな}じ 5の まとまりです。
\end{minipage}\hfill
\begin{minipage}[c]{0.19\linewidth}\centering
\begin{tikzpicture}[x=1mm, y=1mm]
  \foreach \x in {0,3,6,9} {\draw[bkInk, line width=1pt, line cap=round] (\x,0) -- (\x,9);}
  \draw[bkInk, line width=1pt, line cap=round] (-1.5,1.5) -- (10.5,7.5);
  \node[font=\sffamily\bfseries\small, text=bkNavy] at (4.5,-4.5) {5};
  \node[inner sep=0pt] at (18,4.5) {\bkqpsei[7]{5}};
  \node[font=\sffamily\bfseries\small, text=bkNavy] at (18,-4.5) {5};
\end{tikzpicture}
\end{minipage}
\end{column}

\clearpage
%% ---------------- p.179 つかわれているところ ----------------
\begin{butsuri}
ぼうグラフは、 \R{工場}{こう|じょう}や くらしで はたらいて います。
\begin{bkfigure}
\begin{tikzpicture}[x=1mm, y=1mm]
  % 工場のグラフ
  \fill[white, rounded corners=3mm] (0,0) rectangle (72,46);
  \bkqptitle{(38,41.5)}{ロボットが つくった おもちゃ}
  \begin{scope}[shift={(13,8.5)}]
    \bkqpaxes{(0,0)}{52}{7}{3.7}{2}{20}{こ}
    \foreach \v/\c [count=\i] in {120/bkBlue!55,130/bkBlue!55,60/bkRed!75,130/bkBlue!55,120/bkBlue!55} {\bkqpbar{({(\i-0.5)*10.4},0)}{6.4}{\v/20*3.7}{\c}}
    \bkqpnames{10.4}{{\R{月}{げつ}},{\R{火}{か}},{\R{水}{すい}},{\R{木}{もく}},{\R{金}{きん}}}
    \node[draw=bkRed, line width=0.8pt, fill=white, rounded corners=1mm, font=\sffamily\bfseries\fontsize{8.5}{10}\selectfont, text=bkRed, inner sep=1.2pt] at (26,16) {あれっ？};
  \end{scope}
  % スマホ
  \begin{scope}[shift={(76,0)}]
  \fill[white, rounded corners=3mm] (0,0) rectangle (46,46);
  \filldraw[fill=bkInk!85, rounded corners=2.5mm] (11,9) rectangle (35,44);
  \fill[white, rounded corners=1mm] (13,12.5) rectangle (33,40.5);
  \draw[bkInk, line width=0.8pt, rounded corners=0.5mm] (16,32) rectangle (28,37.5); \fill[bkInk] (28,33.8) rectangle (29,35.7);
  \fill[bkGreen] (17,33) rectangle (24.5,36.5);
  \foreach \i/\h in {0/1.5,1/2.5,2/3.5,3/4.5} {\fill[\ifnum\i<3 bkInk\else bkLine\fi] (16.5+\i*2.4,22.5) rectangle ++(1.6,\h);}
  \node[font=\sffamily\bfseries\fontsize{7.5}{8}\selectfont, text=bkInk] at (22,29.6) {でんち};
  \node[font=\sffamily\bfseries\fontsize{7.5}{8}\selectfont, text=bkInk] at (21,18.8) {でんぱ};
  \node[font=\sffamily\bfseries\small] at (23,4.5) {スマホ};
  \end{scope}
\end{tikzpicture}
\end{bkfigure}
\begin{itemize}[label={\color{bkPurple}●}, leftmargin=1.4\zw, itemsep=0pt, topsep=0pt]
\item \R{工場}{こう|じょう}では、 \R{毎日}{まい|にち} つくった \R{数}{かず}を グラフに します。 \R{水曜日}{すい|よう|び}だけ ひくい！ きかいを しらべます。
\item スマホは、 \R{電池}{でん|ち}の のこりや \R{電波}{でん|ぱ}の つよさを、 ぼうの \R{長}{なが}さで \R{見}{み}せます。
\item \R{天気}{てん|き}よほうの こうすいかくりつや、 せんきょの ニュースにも グラフが \R{出}{で}ます。
\end{itemize}
\end{butsuri}

\bkfukidashi{delta}{\R{工場}{こう|じょう}の \R{人}{ひと}は、 \R{表}{ひょう}と グラフで きかいの ちょうしを \R{見}{み}まもって いるよ。 \R{表}{ひょう}は、 コンピュータの \R{計算}{けい|さん}にも つかわれるんだ。}

\begin{bkfigure}
\begin{tikzpicture}[x=1mm, y=1mm]
  \draw[bkRoute5, line width=3pt, line cap=round] (10,0) -- (124,0);
  \node[anchor=west, inner sep=0pt] at (0,12) {\bkrouteicon[1.4em]{5}\hspace{0.3em}{\sffamily\bfseries\small\color{bkRoute5!80!black}\R{表}{ひょう}の \R{路線}{ろ|せん}}};
  \foreach \x/\c/\t/\s in {10/bkVol1/{1-13}/{\R{九九}{く|く}の \R{表}{ひょう}},48/bkVol1/{1-17}/{\R{表}{ひょう}と グラフ},86/bkVol3/{3-6}/{れんりつ},124/bkVol7/{7-1}/{\R{行列}{ぎょう|れつ}}} {
    \fill[white, draw=\c, line width=1.6pt] (\x,0) circle (2.6);
    \node[anchor=base, font=\sffamily\bfseries\footnotesize, text=\c!85!black] at (\x,-6.8) {\t};
    \node[anchor=base, font=\sffamily\footnotesize, text=bkInk] at (\x,-13.6) {\s};}
  \node[fill=bkTheme, text=white, rounded corners=1mm, font=\sffamily\bfseries\fontsize{8}{9}\selectfont, inner sep=1.5pt] at (48,5.5) {いま ここ};
\end{tikzpicture}
\end{bkfigure}

\clearpage
%% ---------------- p.180〜181 もんだい ----------------
\bkendproblems

\bkproblemgroup{きほん}{1}

\bkproblem \R{正}{せい}の \R{字}{じ}を \R{数}{かず}に しましょう。\par\vspace{2pt}
\noindent\hspace*{1\zw}（1）\bkqpsei{11}\enspace → \bkkakikomi\,\R{人}{にん}\hspace{3\zw}（2）\bkqpsei{8}\enspace → \bkkakikomi\,\R{人}{にん}
\bkpans{（1）11\R{人}{にん}\quad（2）8\R{人}{にん}}

\bkproblem
\begin{minipage}[t]{0.46\linewidth}
（1）いちばん \R{多}{おお}い のりものは？\par
\bkkakikomi[26mm]\par
（2）ひこうきと ふねの ちがいは \R{何人}{なん|にん}？
\end{minipage}\hfill
\begin{minipage}[t]{0.44\linewidth}\centering\vspace{0pt}
\begin{tikzpicture}[x=1mm, y=1mm]
  \bkqptitle{(26,33.5)}{すきな のりもの}
  \bkqpgraph{(0,0)}{52}{10}{2.8}{2}{1}{\R{人}{にん}}{13}{9/ひこうき,{8}/{\R{電車}{でん|しゃ}},5/バス,3/ふね}
\end{tikzpicture}
\end{minipage}
\bkshikiran
\bkpans{（1）ひこうき\quad（2）9 － 3 ＝ 6\quad 6\R{人}{にん}}

\bkproblem
\begin{minipage}[t]{0.46\linewidth}
（1）1\R{目}{め}もりは \R{何人}{なん|にん}？\par
\hspace*{1\zw}\bkkakikomi\,\R{人}{にん}\par
（2）あきは \R{何人}{なん|にん}？\enspace\bkkakikomi\,\R{人}{にん}
\end{minipage}\hfill
\begin{minipage}[t]{0.44\linewidth}\centering\vspace{0pt}
\begin{tikzpicture}[x=1mm, y=1mm]
  \bkqptitle{(26,30)}{すきな きせつ}
  \bkqpgraph{(0,0)}{52}{6}{4.2}{2}{10}{\R{人}{にん}}{13}{40/はる,60/なつ,50/あき,30/ふゆ}
\end{tikzpicture}
\end{minipage}
\bkpans{（1）10\R{人}{にん}\quad（2）50\R{人}{にん}}

\bkproblem
\begin{minipage}[t]{0.46\linewidth}
\R{見}{み}つけた \R{虫}{むし}の \R{数}{かず}を、 ぼうグラフに かきましょう。
\par\vspace{3pt}
\begin{bktable}{colspec={Q[c,wd=11mm] Q[c,wd=15mm] Q[c,wd=15mm] Q[c,wd=15mm]}, colsep=1.5pt, rowsep=2pt, cells={font=\small}}
アリ & バッタ & チョウ & トンボ \\
14 & 10 & 6 & 4 \\
\end{bktable}
\end{minipage}\hfill
\begin{minipage}[t]{0.44\linewidth}\centering\vspace{0pt}
\begin{tikzpicture}[x=1mm, y=1mm]
  \bkqpaxes{(0,0)}{52}{8}{3.3}{2}{2}{ひき}
  \bkqpnames{13}{アリ,バッタ,チョウ,トンボ}
\end{tikzpicture}
\end{minipage}
\bkpans{ぼうの \R{高}{たか}さ：アリ 7、 バッタ 5、 チョウ 3、 トンボ 2（\R{目}{め}もり）}

\clearpage
\bkproblemgroup{チャレンジ}{1}

\bkproblem
\begin{minipage}[t]{0.46\linewidth}
（1）いちばん \R{少}{すく}ない \R{日}{ひ}は？\par
\bkkakikomi[14mm]\par
（2）\R{月曜日}{げつ|よう|び}と \R{水曜日}{すい|よう|び}の ちがいは？
\end{minipage}\hfill
\begin{minipage}[t]{0.44\linewidth}\centering\vspace{0pt}
\begin{tikzpicture}[x=1mm, y=1mm]
  \bkqptitle{(26,33.5)}{うれた パンの \R{数}{かず}}
  \bkqpgraph{(0,0)}{52}{10}{2.8}{5}{10}{こ}{10.4}{70/{\R{月}{げつ}},50/{\R{火}{か}},90/{\R{水}{すい}},60/{\R{木}{もく}},80/{\R{金}{きん}}}
\end{tikzpicture}
\end{minipage}
\bkshikiran
\bkpans{（1）\R{火曜日}{か|よう|び}\quad（2）90 － 70 ＝ 20\quad 20こ}

\bkproblem すきな \R{本}{ほん}を しらべました。 10\R{目}{め}もりの グラフに かくとき、 1\R{目}{め}もりは 1\R{人}{にん}・2\R{人}{にん}・5\R{人}{にん}の どれに する？
\par\vspace{3pt}
\noindent\begin{minipage}[c]{0.64\linewidth}
\begin{bktable}{colspec={Q[c,wd=14mm] Q[c,wd=13mm] Q[c,wd=13mm] Q[c,wd=12mm] Q[c,wd=13mm]}, colsep=1.5pt, rowsep=2pt, cell{1-2}{1}={bg=bkTheme!12, font=\sffamily\bfseries\small}, cells={font=\small}}
\R{本}{ほん} & まんが & ずかん & \R{絵本}{え|ほん} & クイズ \\
\R{人数}{にん|ずう} & 45 & 20 & 35 & 10 \\
\end{bktable}
\end{minipage}\hfill
\begin{minipage}[c]{0.34\linewidth}\raggedleft
1\R{目}{め}もり\enspace\bkkakikomi\,\R{人}{にん}
\end{minipage}
\bkpans{5\R{人}{にん}（45\R{人}{にん}は 9\R{目}{め}もり。 1\R{人}{にん}や 2\R{人}{にん}では 10\R{目}{め}もりに \R{入}{はい}らない）}

\begin{furikaeri}
\begin{checklist}
\item しらべて かぞえ、 \R{表}{ひょう}に まとめられる。\hfill{\footnotesize\sffamily\color{bkTheme}→\pageref{bk:1-17:hyou}ページ}
\item ぼうグラフが かけて、 よめる。\hfill{\footnotesize\sffamily\color{bkTheme}→\pageref{bk:1-17:bou}ページ}
\item グラフの 1\R{目}{め}もりを たしかめられる。\hfill{\footnotesize\sffamily\color{bkTheme}→\pageref{bk:1-17:memori}ページ}
\end{checklist}
チェックが つかなかったら、 →の ページに もどろう。
\end{furikaeri}

\bkfukidashi{sora}{おうちの \R{人}{ひと}の すきな たべものを しらべて、 ぼうグラフに して みよう！}
\clearpage
